\section{Numerical methods}
\label{sec:numerical-framework:discretization}
% Configuration provenance (WSL):
% /home/mgnadlinger/Zhang_DNS_ZPGFP/kwSST_Prt0p9/M6Tw025/system/fvSchemes
% /home/mgnadlinger/Zhang_DNS_ZPGFP/kwSST_Prt0p9/M6Tw025/system/fvSolution
% Solver source: /home/mgnadlinger/OpenFOAM/mgnadlinger-v2412/modules/hysa/
% src/finiteVolume/compactSnGradScheme/compactSnGrad.{H,C}
% src/finiteVolume/fluxSchemes/ausmDV/ausmDVFluxScheme.{H,C}

The equations are discretized by a cell-centred finite-volume method.
Fluxes across cell faces determine the discrete conservation balances.
Table~\ref{tab:numerical-framework:schemes} summarizes the numerical setup, which
unless stated otherwise, is the same numerical method used throughout this
work, with the turbulence-equation entries adapted to the selected model.
Pseudo-time Courant limits and other convergence controls are adjusted by
case. The complete settings are supplied with the simulation setups in the
\hyperref[chap:data-avail]{Data Availability statement}.

\begin{table}[htbp]
    \centering
    \small
    \caption{Numerical methods in the representative M6Tw025 configuration.}
    \label{tab:numerical-framework:schemes}
    \begin{tabular}{@{}p{0.29\linewidth}p{0.65\linewidth}@{}}
        \toprule
        Operation & Selected method \\
        \midrule
        Convective flux & AUSMDV2\(^+\)-based scheme;
            dictionary selection \texttt{AUSMDV}
            (Section~\ref{subsec:ext:AUSM}) \\
        Face-state reconstruction & Weighted van Leer limiter
            (\texttt{wVanLeer}) for density, velocity and temperature \\
        Cell gradients & Limited Gauss gradients:
            \texttt{cellLimited Gauss linear 1} \\
        Momentum and energy diffusion & \texttt{Gauss linear compact}:
            linear coefficient interpolation and HiSA's compact
            face-normal gradient \\
        Other diffusion operators & \texttt{Gauss linear corrected};
            general face-normal gradient \texttt{corrected} \\
        Turbulence convection & Bounded first-order upwind for
            turbulence properties \\
        \bottomrule
    \end{tabular}
\end{table}

\subsection{Convective and diffusive face operations}

Convective fluxes use reconstructed states on either side of each face.
The reconstruction and its limiting must be distinguished from the
face-normal gradients used for viscous stress and thermal diffusion.
\par\medskip

For an orthogonal wall-normal stencil, the compact gradient reduces to the
slope between neighbouring cell-centre values. The coefficient multiplying
that gradient is interpolated to the face with distance weighting
\begin{equation}
    D_f\phi=\frac{\phi_R-\phi_L}{h},
    \qquad
    \Gamma_f\simeq\frac{b\Gamma_L+a\Gamma_R}{h},
    \qquad
    a=y_f-y_L,\quad b=y_R-y_f,\quad h=a+b.
    \label{eq:numerical-framework:internal-face-operators}
\end{equation}
Here \(y_L\) and \(y_R\) are the neighbouring cell-centre positions, \(y_f\)
is the face position, and \(h=y_R-y_L\) is the full centre separation.
The field \(\phi\) denotes a transported quantity, such as internal energy,
and \(\Gamma\) its effective diffusion coefficient. For a perfect gas with
constant specific heat, the energy-diffusion term can equivalently be
expressed using temperature and effective thermal conductivity. On non-orthogonal
meshes, the compact scheme adds a correction using interpolated cell
gradients, but in the scope of this work all meshes are 2D, structured, orthogonal grids.
\par\medskip

For a face midway between two centres, the secant gradient has a
second-order truncation error. Unequal centre-to-face distances introduce
an additional asymmetry term. Thus, neither the nominal order nor the
wall-function sampling height alone establishes the accuracy of an
internal-face flux. The cold-wall profile analysis in
Section~\ref{subsec:comparative:zpg:mesh-criterion} examines these gradient
and coefficient-interpolation operations and their implications.
